\section{Denominator compatibility and the effective activation loss}
\label{sec:tpc64-compatible-chain}

A terminal numerator becomes useful only when every paper in the chain
uses the same parent denominator.  This section records the exact
compatibility requirements, defines the inflation and effective terminal
losses, and inserts them into the inherited endpoint ledger without
double charging.

\subsection{The literal compatibility audit}

Fix either a declared coefficient class
\(\mathfrak Z_X\) or one distinguished singleton
\(\{z_*(X)\}\).  Every item below must concern that same scope.

\begin{definition}[Compatible realization]
\label{def:tpc64-denominator-compatible-realization}
A realization of the TPC--47/51/54/59/62/63 chain is
\emph{denominator compatible} if all of the following hold
\citep{WangTPC47,WangTPC51,WangTPC54,WangTPC59,WangTPC62,WangTPC63}.
\begin{enumerate}[label=\textup{(D\arabic*)},leftmargin=3.4em]
 \item The TPC--47 physical slice uses one fixed \(h_0\), its literal row
       coefficient \(z_m\), one cofactor sign \(\mu(r)\), and the actual
       coefficient-independent full incidence.  No frozen shift or
       independently selected coefficient vector replaces this slice.
 \item The TPC--51 direct sum retains the same
       \((m,r,j,\gamma)\)-atoms and their raw counting-measure norm.
       A Mellin \(L^1\) envelope, a quotient-frame norm, or a
       representation-dependent layer norm is not substituted for the
       physical parent.
 \item The TPC--54 capacity weights are exactly the row masses of the
       parent used here.  After a populated repair they are
       \(\widetilde P_m\), so both \(\delta_\eta\) and \(\kappa_X\) are
       computed with \(\widetilde P_m\), not with source-row cardinality,
       \(P_m\), or a terminal-active submeasure.
 \item The TPC--59 terminal multiplier and support predicate give exactly
       \(T_m\) on the same atoms.  The compatible pre-terminal
       denominator is not replaced by the terminal-support mass \(H_m\).
 \item The TPC--63 provenance test is applied before pair
       rectangularization.  Pair-bad blocks retain their literal active incidence;
       on pair-good blocks every added parent atom and its inflation remain
       in \(\widetilde P_m\).  Terminal-active support is not used
       retroactively as the parent.
 \item In the TPC--62 chain, the form leaving this activation module is
       literally the form entering the represented/native module.  Every
       external block label is retained, or its actual erasure map and
       range-restricted lower bound are inserted as a separate reassembly
       arrow.
\end{enumerate}
\end{definition}

These conditions are identities of objects, not estimates.  When they
hold, the relevant beginning of the chain is
\begin{equation}
 \boxed{
 P(z)
 \xrightarrow[\text{parent completion}]{\widetilde P=P+\Delta}
 \widetilde P(z)
 \xrightarrow[\text{literal terminal map}]{\quad}
 T(z).}
 \label{eq:tpc64-exact-denominator-chain}
\end{equation}
The first arrow is not a lower transfer: it records a possible
denominator enlargement.  The lower constant for the second arrow must
therefore be measured by \(T(z)/\widetilde P(z)\).

\subsection{Exact constants and exponent definitions}

Assume first that the declared class contains at least one vector with
positive literal parent.  Define
\begin{align}
 \mathfrak I_X(\mathfrak Z_X)
 &:=
 \sup_{\substack{z\in\mathfrak Z_X\\\widetilde P(z)>0}}
 \frac{\widetilde P(z)}{P(z)}
 \in[1,+\infty],
 \label{eq:tpc64-compatible-inflation-constant}\\
 c_{\rm term}^{\rm lit}(X;\mathfrak Z_X)
 &:=
 \inf_{\substack{z\in\mathfrak Z_X\\P(z)>0}}
 \frac{T(z)}{P(z)},
 \label{eq:tpc64-literal-terminal-constant}\\
 c_{\rm term}^{\rm eff}(X;\mathfrak Z_X)
 &:=
 \inf_{\substack{z\in\mathfrak Z_X\\
                  \widetilde P(z)>0}}
 \frac{T(z)}{\widetilde P(z)}.
 \label{eq:tpc64-effective-terminal-constant}
\end{align}
In \eqref{eq:tpc64-compatible-inflation-constant}, the quotient is
understood as \(+\infty\) when
\(P(z)=0<\widetilde P(z)\).  Thus a finite compatible inflation
constant automatically excludes a new-only vector from the declared
scope.
For the canonical hybrid repair there is no new-only row, and the
full-space value of \(\mathfrak I_X\) is the maximum row inflation.  For a
distinguished singleton, all three quantities are ratios on that one
vector and must retain the distinguished label.

The exact comparison is
\begin{equation}
 \boxed{
 c_{\rm term}^{\rm eff}(X;\mathfrak Z_X)
 \ge
 \frac{c_{\rm term}^{\rm lit}(X;\mathfrak Z_X)}
      {\mathfrak I_X(\mathfrak Z_X)}}
 \label{eq:tpc64-compatible-activation-comparison}
\end{equation}
whenever the right side is defined with finite
\(\mathfrak I_X\).  Indeed, finiteness gives \(P(z)>0\) whenever
\(\widetilde P(z)>0\), and for every such vector
\[
 \frac{T(z)}{\widetilde P(z)}
 =
 \frac{T(z)}{P(z)}
 \left(\frac{\widetilde P(z)}{P(z)}\right)^{-1}
 \ge
 \frac{c_{\rm term}^{\rm lit}(X;\mathfrak Z_X)}
      {\mathfrak I_X(\mathfrak Z_X)}.
\]
Taking the infimum proves
\eqref{eq:tpc64-compatible-activation-comparison}.  This is sharp
already in a one-row model.  A direct estimate relative to
\(\widetilde P\), such as
\cref{thm:tpc64-threshold-capacity}, can be stronger.

For an asymptotic family of scopes, define the optimal fixed-power losses
\begin{align}
 \lambda_{\rm infl}
 &:=
 \limsup_{X\to\infty}
 \frac{\log \mathfrak I_X(\mathfrak Z_X)}{\log X},
 \label{eq:tpc64-lambda-infl-definition}\\
 \lambda_{\rm term}^{\rm eff}
 &:=
 \limsup_{X\to\infty}
 \frac{\log^+\!\bigl(
       1/c_{\rm term}^{\rm eff}(X;\mathfrak Z_X)\bigr)}
      {\log X},
 \label{eq:tpc64-lambda-term-eff-definition}
\end{align}
provided the corresponding constant is positive eventually.  If
\(\mathfrak I_X=+\infty\) or
\(c_{\rm term}^{\rm eff}=0\) infinitely often, assign the relevant loss
\(+\infty\).  Here \(\log^+u:=\max\{\log u,0\}\).  Define
\(\lambda_{\rm term}^{\rm lit}\) analogously from
\(c_{\rm term}^{\rm lit}\).

\begin{proposition}[Effective loss from the two honest routes]
\label{prop:tpc64-two-effective-loss-routes}
In one fixed nonempty quantifier scope.  In route \textup{(ii)} this scope
is a normalized or capacity-restricted slice and need not be homogeneous:
\begin{enumerate}[label=\textup{(\roman*)},leftmargin=3em]
 \item if
       \(\mathfrak I_X\le X^{\lambda_{\rm infl}+o(1)}\) and
       \(c_{\rm term}^{\rm lit}
          \ge X^{-\lambda_{\rm term}^{\rm lit}+o(1)}\),
       then
       \begin{equation}
        c_{\rm term}^{\rm eff}
        \ge
        X^{-(\lambda_{\rm term}^{\rm lit}
             +\lambda_{\rm infl})+o(1)};
        \label{eq:tpc64-inflated-terminal-exponent}
       \end{equation}
 \item if common \(C_X>0\) and \(0<\kappa_X\le1\) exist such that every
       \(z\in\mathfrak Z_X\) satisfies
       \[
        |z_m|\le C_X
        \quad(m\in\mathcal M_X^+),
        \qquad
        \widetilde P(z)
        \ge\kappa_XC_X^2\widetilde W_X>0,
       \]
       and the threshold-capacity inputs give
       \begin{equation}
        \eta_X
        \left(
        1-\frac{1-\delta_{\eta_X}}{\kappa_X}
        \right)_+
        \ge X^{-\lambda_{\rm cap}+o(1)},
        \label{eq:tpc64-capacity-exponent-input}
       \end{equation}
       on the same declared class, then
       \begin{equation}
        c_{\rm term}^{\rm eff}
        \ge X^{-\lambda_{\rm cap}+o(1)}.
        \label{eq:tpc64-capacity-effective-exponent}
       \end{equation}
\end{enumerate}
If both routes apply, the larger lower constant may be used.  In
particular,
\begin{equation}
 \lambda_{\rm term}^{\rm eff}
 \le
 \min\{
 \lambda_{\rm term}^{\rm lit}+\lambda_{\rm infl},
 \lambda_{\rm cap}
 \}.
 \label{eq:tpc64-best-effective-loss}
\end{equation}
No equality is asserted: a direct terminal estimate may be better than
either certificate.
\end{proposition}

\begin{proof}
The first conclusion is
\eqref{eq:tpc64-compatible-activation-comparison} in exponent form.  The
second is \cref{thm:tpc64-threshold-capacity}.  Taking the better of two
proved lower bounds gives \eqref{eq:tpc64-best-effective-loss}.
\end{proof}

If \(\mathfrak Z_X=\{z_*(X)\}\), every constant and exponent above is
distinguished.  A bound proved there cannot be inserted into the
full-space versions of
\eqref{eq:tpc64-compatible-inflation-constant}--%
\eqref{eq:tpc64-effective-terminal-constant}.

\subsection{The compatible physical chain}

Let \(\mathcal Q_{\rm in}\), \(\mathcal Q_A\),
\(\mathcal Q_F\), and \(\mathcal Q_{\rm out}\) denote, respectively, the
declared upstream input, represented native output, full
represented-plus-missing output, and downstream reconnected output.  The
terminal form \(T\) and repaired parent \(\widetilde P\) are inserted as
two additional named nodes.

\begin{theorem}[Exact five-arrow compatible chain]
\label{thm:tpc64-five-arrow-compatible-chain}
Suppose all six forms are nonnegative quadratic forms evaluated on the same
initial \(z\) in the same scope.  Suppose also that
\[
 c_{\rm pre},c_{\rm term},c_{\rm post},c_{\rm reasm},c_{\rm out}\ge0
\]
and that
\begin{align}
 \widetilde P(z)
 &\ge c_{\rm pre}(X)\mathcal Q_{\rm in}(z),
 \label{eq:tpc64-pre-parent-arrow}\\
 T(z)
 &\ge c_{\rm term}(X)\widetilde P(z),
 \label{eq:tpc64-terminal-arrow}\\
 \mathcal Q_A(z)
 &\ge c_{\rm post}(X)T(z),
 \label{eq:tpc64-post-terminal-arrow}\\
 \mathcal Q_F(z)
 &\ge c_{\rm reasm}(X)\mathcal Q_A(z),
 \label{eq:tpc64-reassembly-arrow}\\
 \mathcal Q_{\rm out}(z)
 &\ge c_{\rm out}(X)\mathcal Q_F(z).
 \label{eq:tpc64-downstream-arrow}
\end{align}
Every output carrier and normalization in one line is required to be the
literal input carrier and normalization in the next.  Then
\begin{equation}
 \boxed{
 \mathcal Q_{\rm out}(z)
 \ge
 c_{\rm out}c_{\rm reasm}c_{\rm post}
 c_{\rm term}c_{\rm pre}\,
 \mathcal Q_{\rm in}(z).}
 \label{eq:tpc64-five-arrow-product}
\end{equation}
The statement is uniform only if all five inequalities and all their
remainders are uniform on the same \(\mathfrak Z_X\).
\end{theorem}

\begin{proof}
Apply the five inequalities successively to the named intermediate forms
generated by the same \(z\).  No supremum, infimum, selector, shift, or
coefficient vector changes during the substitutions.
\end{proof}

The exact TPC--47/51 dictionary may make
\(c_{\rm pre}=1\) on a literally identical carrier.  TPC--54 selection
or TPC--59 occupancy costs belong to \(c_{\rm pre}\) only if they compare
\(\mathcal Q_{\rm in}\) with the precise repaired parent used in
\eqref{eq:tpc64-terminal-arrow}.  TPC--62 reassembly begins only at
\eqref{eq:tpc64-reassembly-arrow}; it cannot repair a zero terminal
constant.

\subsection{No-double-charge ledger and stop rule}

Write
\begin{equation}
 \Lambda_{\rm pre}
 :=\lambda_{\rm sel}+\lambda_{\rm occ}+\lambda_{\rm blk}.
 \label{eq:tpc64-pre-terminal-ledger}
\end{equation}
When each symbol corresponds to a genuinely distinct arrow in
\cref{thm:tpc64-five-arrow-compatible-chain}, the inherited physical
ledger is
\begin{equation}
 \boxed{
 \begin{aligned}
 \Lambda_{\rm tot}^{(64)}
 :={}&
 \lambda_{\rm sel}+\lambda_{\rm occ}+\lambda_{\rm blk}
 +\lambda_{\rm term}^{\rm eff}
 +\lambda_{\rm sing}\\
 &+\lambda_H+\lambda_{\rm aff}
 +\lambda_{\rm reasm}
 +\lambda_{\rm bdry}+\lambda_{\rm tail}.
 \end{aligned}}
 \label{eq:tpc64-complete-loss-ledger}
\end{equation}
The strict endpoint test is
\begin{equation}
 \boxed{\Lambda_{\rm tot}^{(64)}<\frac1{400}.}
 \label{eq:tpc64-strict-stop-rule}
\end{equation}

The following no-double-charge rules are part of the formula.
\begin{enumerate}[label=\textup{(\alph*)},leftmargin=2.8em]
 \item If
       \(\lambda_{\rm term}^{\rm eff}\) is obtained from
       \(\lambda_{\rm term}^{\rm lit}+\lambda_{\rm infl}\), only
       \(\lambda_{\rm term}^{\rm eff}\) enters
       \eqref{eq:tpc64-complete-loss-ledger}; the two diagnostic inputs are
       not added again.
 \item If the capacity theorem proves the terminal arrow directly relative
       to \(\widetilde P\), then
       \(\eta_X,\delta_{\eta_X},\kappa_X\), and
       \(\lambda_{\rm cap}\) are ingredients of the one term
       \(\lambda_{\rm term}^{\rm eff}\), not four further arrows.
 \item If selection, occupancy, or block losses were already used to prove
       the same direct ratio \(T/\widetilde P\), they must be removed from
       \(\Lambda_{\rm pre}\).  Conversely, if they compare an earlier
       \(\mathcal Q_{\rm in}\) with \(\widetilde P\), they remain separate.
 \item A represented-reference exponent used only to certify
       \(\lambda_{\rm reasm}\) is not charged again if its physical arrow
       is already included in \(\lambda_H+\lambda_{\rm aff}\).
 \item Exact set identities, faithful hypergraph retention, literal
       coordinate permutations or unitary relabelings, and \(O(\log X)\)
       diagnostic labels have zero fixed-power cost.  A more general
       coordinate change has zero such cost only after a uniform or
       \(X^{o(1)}\) condition-number bound is proved; otherwise its loss
       must enter the ledger.  Erasing diagnostic labels is not a
       coordinate change and belongs to \(\lambda_{\rm reasm}\).
\end{enumerate}

If a mandatory constant in
\cref{thm:tpc64-five-arrow-compatible-chain} is zero infinitely often,
the matching exponent route stops.  If a necessary proved loss is at
least \(1/400\), the exponent-only endpoint route also stops, including at
equality.  A separate constant-sensitive endpoint argument at equality
would lie outside this strict exponent ledger and would not satisfy
\eqref{eq:tpc64-strict-stop-rule}.  A distinguished success, an
averaged-shift success, or a different-denominator success does not evade
either stop rule.
